% ============================================================
% Chapter: Reliability and Probabilistic Assessment
% Assessed MLOs: MLO-3
% Excellent-grade rubric criteria:
% - Application of quantified uncertainties (5%): key uncertainties are defined,
%   parameterised, and propagated in a way that is relevant to design decisions.
% - Concept design probabilistic assessment quality (5%): component and system
%   reliability results are technically sound, interpretable, and used to identify
%   design improvements or risk drivers.
% ============================================================

\section{Reliability and Probabilistic Assessment}
\pagecheck{5--6}

% Rubric Checklist: Application of quantified uncertainties (5%); Concept design probabilistic assessment quality (5%).
\instructionplaceholder{Uncertainty definition}{Identify the uncertainties that materially affect the chosen concept. Distinguish load, resistance, model, and operational uncertainties, and state which ones are included in the probabilistic model and which ones are only discussed qualitatively.}

\subsection{Probabilistic model formulation}
\instructionplaceholder{Required evidence}{Define the random variables, their distributions or bounds, and the reasoning used to select them. State the key limit state or performance function, the reliability method applied, and the assumptions needed to make the analysis tractable.}

\subsection{Component reliability assessment}
\instructionplaceholder{Required evidence}{Select at least one critical component or limit state and show how its failure probability or reliability index is evaluated. Explain why this component is critical for the concept and what the result implies for design robustness.}

\subsection{System-level risk assessment}
\instructionplaceholder{Required evidence}{Extend the analysis to a simple system view. Use a fault tree, event tree, or other clear logic to show how component-level failures combine into broader system risk. The goal is not complexity; the goal is to show that the concept has been assessed beyond a single isolated check.}

\subsection{Interpretation and recommendations}
\instructionplaceholder{Required evidence}{Interpret the probabilistic results in engineering terms. Identify the variables or assumptions that dominate the result, explain what safety margin or vulnerability is revealed, and state what modifications would most effectively improve reliability.}
